\section{Introduction}\label{sec:intro}

Nine languages are spoken on the shores of the Baltic Sea. They belong to four families---Finnic, Baltic, Slavic and Germanic---that have traded, fished, fought and intermarried across the same water for a thousand years, and yet they mostly cannot understand one another. The sea itself has no name of its own. Germans, Swedes and Danes call it the East Sea (\emph{Ostsee}, \emph{Östersjön}, \emph{Østersøen}); Finns, calquing the Swedish, also call it the East Sea (\emph{Itämeri}) although it lies to their west; Estonians call it the West Sea (\emph{Läänemeri}); Latvians, Lithuanians, Poles and Russians call it the Baltic after a Latin coinage of Adam of Bremen around 1075, whose etymology is itself uncertain \citep{adam1075,jezierski2016}. Each shore names the sea from where it stands. None names it from the sea.

This paper describes \textbf{Baltmeri}, a constructed language designed to be the sea's own: a language whose every word is decided by a fixed, repeatable procedure applied to the nine coastal languages, so that no nation owns it and every nation can hear itself in it. The name is the procedure's first product: \emph{balt-} is shared by Baltic, Slavic and Germanic speakers, and \emph{meri} is what Finns and Estonians call the sea, and what Slavs almost call it (\emph{more}). We present the language's design goals, its phonology, morphology, syntax and lexicon, the algorithm that generates its vocabulary, an evaluation of how faithfully that algorithm reproduces the language's founding text, and a translation system that makes the language usable today.

\subsection{Why the sea needs a voice}
The motivation is ecological before it is linguistic. The Baltic is shallow, brackish and semi-enclosed; its water takes about thirty years to exchange with the North Sea, and its catchment is home to more than 85 million people in fourteen countries \citep{helcom2023}. It is, in the words of \citet{carstensen2014}, ``the largest anthropogenically induced hypoxic area in the world'': oxygen-depleted bottom water has expanded roughly tenfold in a century, from about 5,000 to more than 60,000~km$^2$. HELCOM's holistic assessments found eutrophication affecting 97\% of the sea's area in 2011--2016 and ``no clear signs of recovery'' in 2016--2021, with only four of fifteen commercial fish stocks in good status \citep{helcom2018,helcom2023}. The eastern Baltic cod, once the region's defining fishery, has been under zero-catch advice since 2019 \citep{ices2019}. The Baltic Proper harbour porpoise numbers around five hundred animals whose range crosses the waters of nine states \citep{amundin2022}. Where the coastal states have acted together, the sea has answered: grey seals, reduced to about 3,000 animals in the early 1980s by organochlorine poisoning, were counted at some 42,000 in 2021 \citep{helcomseal}. \citet{reusch2018} draw the general lesson: the Baltic's trend reversals ``could be achieved only by implementing an international, cooperative governance structure transcending its complex multistate policy setting.''

Institutions for that cooperation exist and are old. The Helsinki Convention dates from 1974; the Council of the Baltic Sea States was founded in 1992 \citep{cbss1992}; the EU Strategy for the Baltic Sea Region of 2009 was the Union's first macro-regional strategy \citep{eusbsr2009}; the Baltic Sea Action Plan was adopted in 2007 and renewed in 2021 \citep{helcombsap2021}. Yet identity has not followed institutions. \citet{gotz2016} maps at least eight overlapping regional frameworks and concludes that ``a Baltic Sea identity has barely evolved beyond the circle of an activist elite''; \citet{lehti2003} describe competing, not converging, region-building projects. Historians have long shown that a sea can be written as a subject rather than as the empty space between nations---\citet{braudel1949} for the Mediterranean, \citet{horden2000,horden2006} for the ``new thalassology'', \citet{north2015} for the Baltic as a ``Nordic Mediterranean''---but the people on its shores have few everyday means of experiencing that subject.

Two further developments make the present moment different. First, the law has begun to recognise ecosystems as persons: the Whanganui River in New Zealand \citeyearpar{teawatupua2017}, the Atrato River in Colombia, and, in 2022, the Mar Menor lagoon in Spain, the first European ecosystem with legal personality, established by a citizens' initiative with 639,824 signatures \citep{spain2022,kauffman2021,stone1972}. In the Nordic region, \citet{thiel2024} has proposed an Embassy of the Baltic Sea as ``the first embassy for more-than-human representation in a transnational context'', and describes the human identity it implies as that of a \emph{bioregional citizen belonging to a watershed}. Second, the environmental humanities have supplied a vocabulary for taking such a subject seriously: Haraway's sympoiesis, Latour's Gaia, Neimanis's \emph{hydrocommons}---``we are all bodies of water''---and Steinberg and Peters's ``wet ontology'' of the sea as voluminous, material and continually reforming \citep{haraway2016,latour2017,neimanis2017,steinberg2015}. Neimanis and Åsberg have written specifically of the Baltic, fathoming the chemical weapons dumped in the Gotland Deep \citep{neimanis2017gotland}.

A person needs a voice. \citet{anderson2006} showed that shared print-languages were the substrate on which national communities imagined themselves into being before they existed politically; ecolinguistics has argued that the stories a language makes easy to tell shape how its speakers treat their environment \citep{haugen1972,fill2001,stibbe2021}. Baltmeri is an attempt to run Anderson's mechanism in reverse and on behalf of the sea: to build, deliberately and transparently, a language in which the Baltic is the centre rather than the edge, in which \emph{meri} ``sea'' and \emph{kala} ``fish'' and \emph{hülje} ``seal'' are everyone's words because they were chosen by everyone's languages, and in which the sentence \emph{Meri er men}---``the sea is ours''---is true by construction. We are all citizens of the Baltic; Baltmeri is what those citizens would speak if the sea did the speaking.

\subsection{Why the procedure must be deterministic}
Constructed auxiliary languages have a long history of committees exercising taste \citep{okrent2009,gobbo2020}. For a language meant to belong to nine nations equally, taste is a liability: any choice a Swede or a Pole makes is a choice a Finn or a Latvian can resent. Baltmeri therefore fixes a \emph{procedure} rather than a vocabulary. Two people who follow the Manual with the same nine dictionaries must arrive at the same word; nothing is chosen by preference, and where the nine languages disagree, a fixed voting rule and a rotating tie-break decide. The procedure treats the four families as equal voters, so that the demographic weight of Russian and German cannot swamp Estonian and Latvian, and it breaks ties by distance from Visby on Gotland, the Hanseatic hub in the middle of the sea, so that the tie-break is the sea's geography and not anyone's politics.

Determinism has a second consequence that turns out to be the language's most distinctive feature: every Baltmeri word can \emph{show its work}. Because the derivation is an algorithm, the translator we describe in \S\ref{sec:engine} can display, for every output word, the nine source forms that voted, which family won, and how the tie was broken. In a region where language has so often been an instrument of domination---German as the Hanseatic lingua franca, Swedish, Russian and German as imperial languages in the Baltic provinces \citep{serzant2025,north2015}---a language that is accountable for every one of its words is a small political proposition of its own.

\subsection{Contributions}
\begin{enumerate}
\item A complete description of Baltmeri: source languages and tie-break geography (\S\ref{sec:sources}), phonology and orthography (\S\ref{sec:phonology}), the seven-rule lexicogenesis procedure (\S\ref{sec:procedure}), morphology (\S\ref{sec:morphology}), syntax (\S\ref{sec:syntax}) and lexicon (\S\ref{sec:lexicon}).
\item An implementation of the procedure as a dependency-free, deterministic engine with a serialisable tie-break state, and a reproducibility study in which the engine re-derives the founding Manual's 58 lexical decrees from their nine source forms (\S\ref{sec:eval}). The engine reproduces 28; we classify the 30 divergences and show that they concentrate in one step of one rule.
\item A translation system in which a large language model performs only source-language analysis and dictionary lookup, the engine derives and inflects every Baltmeri word, and a second model call writes an explanation grounded in the engine's trace (\S\ref{sec:engine}), deployed publicly.
\item An argument, developed in \S\ref{sec:discussion}, that an algorithmically synthesised regional language is a legitimate instrument of more-than-human representation, and a candid account of where the Manual's author overrode the algorithm and what that teaches about hand-made versus generated language.
\end{enumerate}
